\section{Proof of the Fisher information theorem}\label{app:fisher}

We keep the notation of Section~\ref{sec:fisher}. Let
$\sigma_i\in\{-1,1\}$ be the direction of step $i$. Then
$X=\sigma_1+\dots+\sigma_N$, and $J$ counts the $i<N$ with
$\sigma_{i+1}\ne\sigma_i$. Put $Z=\E[J\mid K]$, so $e_N=\Var Z/\Var J$
by Proposition~\ref{prop:eta}.

\begin{lemma}[exact moments]\label{lem:fisher-moments}
For $N\ge1$ and $0<q<1$,
\begin{align*}
 \E X^2&=\frac{Np}q-\frac{\rho(1-\rho^N)}{2q^2},\\
 \Cov(J,X^2)&=-\frac p{q^2}\bigl[Nq(1+\rho^N)-p(1-\rho^N)\bigr],\\
 \E X^4&=\frac{3N^2p^2}{q^2}-\frac{Np(1+10\rho+\rho^2)}{2q^3}
 +\frac{\rho(1+2\rho)(2+\rho)(1-\rho^N)}{2q^4}
 -\frac{3Np\rho^{N+1}}{q^3},
\end{align*}
and $\Var X^2=2N^2p^2V/q^2$, with $V$ as in Theorem~\ref{thm:fisher}.
\end{lemma}

\begin{proof}
We show that each moment, as a sequence in $N$, is a combination of the
sequences $N^i$ and $N^i\rho^N$ with $i\le4$. Two such combinations that agree
at ten values of $N$ agree everywhere, and a script compares them at those
values exactly.

Let $m_N(\phi,\chi)=\E e^{\phi X+\chi J}$. With
$A_{ab}=\Prob(\sigma_2=b\mid\sigma_1=a)\,e^{\phi b+\chi\mathbf 1[a\ne b]}$ and
$v_a=\frac12e^{\phi a}$ we have $m_N=v^\top A^{N-1}\ones$. So
$\sum_{N\ge1}m_Nz^N=U/D$ as formal power series, where
$U(z)=zv^\top\operatorname{adj}(I-zA)\ones$ and $D(z)=\det(I-zA)$ have degree
at most $2$ in $z$. For a partial derivative
$\partial$ of order $m$ in $(\phi,\chi)$, induction and the quotient rule give
$\partial(U/D)=U_\partial/D^{m+1}$ with $\deg U_\partial\le2m+2$, as each
derivative multiplies the numerator by $D$ or a derivative of $D$. At
$\phi=\chi=0$ the matrix $A$ has trace $1+\rho$ and determinant $\rho$, so
$D(z)=(1-z)(1-\rho z)$. Let $q\ne\frac12$, so its roots $1$ and
$1/\rho$ are distinct. By partial
fractions $U_\partial/D^{m+1}$ is a constant plus a combination of
$(1-z)^{-j}$ and $(1-\rho z)^{-j}$ with $j\le m+1$, whose coefficients of
$z^N$ are $\binom{N+j-1}{j-1}$ and $\binom{N+j-1}{j-1}\rho^N$. So for $N\ge1$
the sequence $\partial m_N$ at $0$ lies in
$\mathcal S_m=\operatorname{span}\{N^i,N^i\rho^N:i\le m\}$. A sequence in
$\mathcal S_m$ satisfies the recurrence with characteristic polynomial
$(t-1)^{m+1}(t-\rho)^{m+1}$, so it is $0$ if it vanishes at
$N=1,\dots,2m+2$.

The moments $\E X^2$, $\E[JX^2]$ and $\E X^4$ are derivatives of $m_N$ of
order at most $4$ at $0$, so they lie in $\mathcal S_4$. The claimed
formulas for $\E X^2$, $\E X^4$ and $\E[JX^2]=\Cov(J,X^2)+(N-1)q\,\E X^2$
lie in $\mathcal S_2$. So it is enough to check them at $N=1,\dots,10$. The
script \texttt{paperC/verify\_sticky\_fisher.py} computes the three moments exactly for
$N\le12$ by a dynamic program with $q$ as a symbol, and every difference is
$0$. For fixed $N$ both sides
are rational in $q$ with no pole at $\frac12$, so the case $q=\frac12$
follows. Finally, with $R$ standing for $\rho^N$ and $\rho R$ for
$\rho^{N+1}$, the claim $\E X^4-(\E X^2)^2=2N^2p^2V/q^2$ is an identity of
rational functions of $N$, $q$ and $R$, which the same script checks with
sympy.
\end{proof}

\begin{proof}[Proof of the closed form in Theorem~\ref{thm:fisher}]
For $N\ge2$, $X^2$ is $N^2$ on the constant paths and $(N-2)^2$ on the
paths that switch only at the last step, so $\Var X^2>0$ and $V>0$ by
Lemma~\ref{lem:fisher-moments}, which also gives
$\Cov(J,X^2)=-\frac{Np}q[1+\rho^N-p(1-\rho^N)/T]$, so
$\Cov(J,X^2)^2/\Var X^2$ is the squared bracket over $2V$. We divide by
$\Var J=(N-1)pq$.
\end{proof}

\begin{lemma}[convexity]\label{lem:fisher-convex}
Let $f$ and $g$ be convex and differentiable on $[t-\delta,t+\delta]$, with
$0\le g-f\le\eta$ and $g''\le\kappa$ there. Then
$|f'(t)-g'(t)|\le\eta/\delta+\kappa\delta/2$.
\end{lemma}

\begin{proof}
By convexity and Taylor's theorem,
$f'(t)\le(f(t+\delta)-f(t))/\delta\le(g(t+\delta)-g(t)+\eta)/\delta
\le g'(t)+\kappa\delta/2+\eta/\delta$, and the lower bound is the same with
$t-\delta$.
\end{proof}

\begin{lemma}[a Bessel ratio]\label{lem:fisher-bessel}
For $t>0$ let $R(t)=I_1(t)/I_0(t)$. Then
$(\sqrt{1+t^2}-1)/t\le R(t)<1$. In particular $0<1-R(t)\le1/t$.
\end{lemma}

\begin{proof}
Since $I_0'=I_1$ and $I_1'(t)=I_0(t)-I_1(t)/t$, we have $R'=1-R/t-R^2$, and
$R(t)=t/2+O(t^3)$ at $0$. At a first point $t_0$ with $R(t_0)=1$ we would
have $R'(t_0)\ge0$, but the equation gives $-1/t_0$. So $R<1$. Put
$s=\sqrt{1+t^2}$ and $\ell=(s-1)/t$. Then $1-\ell/t-\ell^2=\ell/t=1/(s+1)$ and
$\ell'=1/(s(s+1))$, so $\ell'\le1-\ell/t-\ell^2$. For $d=R-\ell$ this gives
$d'\ge-d/t-(R^2-\ell^2)=-d(1/t+R+\ell)$. So $\Theta d$ is nondecreasing, where
$\Theta(t)=t\exp\int_0^t(R+\ell)$. As $\ell=t/2+O(t^3)$, we have
$d=O(t^3)$ and $\Theta d\to0$ at $0$. Hence $d\ge0$. Finally
$1-\ell=(t+1-s)/t\le1/t$.
\end{proof}

\begin{lemma}[moments of the endpoint]\label{lem:fisher-W}
\begin{enumerate}[(a)]
\item For each $m\ge1$ there is $C_m$ with $\E W^{2m}\le C_mT^{-m}$ for all
$0<q<1$ and $T\ge1$.
\item If $0<q\le\frac12$, then $\E W^4\le3p^2/T^2+9/(2T^4)$.
\end{enumerate}
\end{lemma}

\begin{proof}
(a) Add a step $\sigma_{N+1}$ to the walk. Since
$\E[\sigma_{i+1}\mid\sigma_{\le i}]=\rho\sigma_i$, the
$\xi_{i+1}=\sigma_{i+1}-\rho\sigma_i$ are martingale differences. Summing
$\sigma_{i+1}-\sigma_i=-2q\sigma_i+\xi_{i+1}$ over $i\le N$ gives
$X=(\sigma_1-\sigma_{N+1}+\Xi)/(2q)$ with $\Xi=\sum_{i=2}^{N+1}\xi_i$. Also
$\xi_{i+1}^2$ is $4q^2$ if step $i+1$ repeats step $i$ and $4p^2\le4$
otherwise. So $\sum\xi_i^2\le4q^2N+4J'$, where $J'\sim\operatorname{Bin}(N,q)$
counts the switches among the $N+1$ steps. Burkholder's
inequality~\cite[Ch.~2]{hallheyde1980} gives
$\E\,\Xi^{2m}\le c_m\E(4q^2N+4J')^m\le c_m8^m((q^2N)^m+\E J'^m)$, and
$q^2N\le T$. The falling factorial moments of $J'$ are at most $T^j$, and
$J'^m$ is a combination of the falling factorials of order $j\le m$ with
nonnegative coefficients. So
$\E J'^m\le b_mT^m$ for $T\ge1$. As $|\sigma_1-\sigma_{N+1}|\le2$,
we get $\E X^{2m}\le(2q)^{-2m}2^{2m-1}(2^{2m}+\E\,\Xi^{2m})\le C_mq^{-2m}T^m$,
and $\E W^{2m}=N^{-2m}\E X^{2m}\le C_mT^{-m}$.

(b) We divide the formula for $\E X^4$ by $N^4$. As $0\le\rho<1$, the two terms
of order $T^{-3}$ are at most $0$ and the one of order $T^{-4}$ is at most
$9/(2T^4)$.
\end{proof}

\begin{proof}[Proof of the bound in Theorem~\ref{thm:fisher}]
Let $\mathcal P$ be the orthogonal projection in $L^2$ onto the span of $1$
and $X^2$. Since $X^2$ is a function of $K$, $\Cov(Z,X^2)=\Cov(J,X^2)$. So
$\Var Z=\Cov(J,X^2)^2/\Var X^2+\E(Z-\mathcal PZ)^2$, and
$e_N-R^2_N=\E(Z-\mathcal PZ)^2/\Var J\ge0$. Also
$\E(Z-\mathcal PZ)^2\le\E(Z-G)^2$ for $G=y(1-W^2/2)+\frac12$. From now on
$0<q\le\frac12$ and $1\le T\le N^{1/4}$, and $C$ is an absolute constant that
may change at each use. Then $p\ge\frac12$, $T\le y\le2T$ and
$\Var J=(N-1)pq\ge T/4$. It suffices to show
$\E(Z-G)^2\le C(T^{-2}+T^4/N)$.

The plan is as follows. By Proposition~\ref{prop:eta}, $Z$ is the logarithmic
derivative of the polynomial $S_{k,N-k}$ at $u$. Paper~A gives a Bessel
function that is close to this polynomial. Closeness of two functions does not
give closeness of their derivatives in general, but it does for convex
functions, and this is Step~1. Step~2 computes the logarithmic derivative of
the Bessel function and finds $y\sqrt{1-W^2}+\frac12$ up to a small error.
Step~3 replaces the square root by $1-W^2/2$. Steps~4 and~5 add up the errors
in mean square.

For $0<k<N$ put $s_k=\sqrt{k(N-k)}$, $c_k=N/(2s_k)$ and $w_k=2s_ku$. On
$K=k$ we have $s_k^2=N^2(1-W^2)/4$, so $c_k=(1-W^2)^{-1/2}\ge1$,
$w_k=y\sqrt{1-W^2}$ and $c_kw_k=y$. Let $\mathcal I_k(z)=I_0(2z)+c_kI_1(2z)$
and $\widetilde S_k(v)=2v\,\mathcal I_k(s_kv)$, which approximates
$S_{k,N-k}(v)$. Here $I_0$ and $I_1$ are the modified Bessel functions. By
the uniform Bessel estimate for every bin in~\cite{paperA} (the theorem
``uniform estimate for every bin'' there, with $a=k$ and $b=N-k$),
\begin{equation}\label{eq:fisher-bessel-bound}
 0\le1-\frac{S_{k,N-k}(v)}{\widetilde S_k(v)}\le\frac{Nv^2}2+v
 \qquad\text{for all }v>0 .
\end{equation}
Let $\Lambda_k(t)=\log S_{k,N-k}(e^t)$ and
$\beta_k(t)=\log\widetilde S_k(e^t)$. By Proposition~\ref{prop:eta},
$Z=\Lambda_K'(\log u)$ on $0<K<N$.

\emph{Step 1.} Both $\Lambda_k$ and $\beta_k$ are convex, as logarithms of
series in $e^t$ with nonnegative coefficients. Let $|t-\log
u|\le1$. Then $v=e^t\le eu$ and $y\le2N^{1/4}$ give
$Nv^2/2+v\le(e^2y^2/2+ey)/N\le(2e^2+2e)/\sqrt N\le\frac12$ for $N\ge1700$. As
$-\log(1-\varepsilon)\le2\varepsilon$ for $0\le\varepsilon\le\frac12$,
\eqref{eq:fisher-bessel-bound} gives $0\le\beta_k-\Lambda_k\le\eta$ there,
with $\eta=(e^2y^2+2ey)/N$. Next, $\beta_k(t)-t-\log2$ is the logarithm of
$\mathcal I_k(s_ke^t)=\sum_na_ne^{nt}$ with $a_n\ge0$. So $\beta_k''$ is the
variance of $n$ under the weights $a_ne^{nt}$, at most the mean of $n^2$. As
$(z\frac d{dz})^2$ maps $I_0(2z)$ to $4z^2I_0(2z)$ and $I_1(2z)$ to
$(4z^2+1)I_1(2z)$, this mean is at most $4z^2+1\le e^2y^2+1=\kappa$, since
$z=s_ke^t\le ey/2$. For $N\ge4$,
$\delta=\sqrt{2\eta/\kappa}\le1$ and Lemma~\ref{lem:fisher-convex} gives
\begin{equation}\label{eq:fisher-step1}
 \bigl|\Lambda_k'(\log u)-\beta_k'(\log u)\bigr|
 \le\sqrt{2\eta\kappa}
 \le\frac{2e^2y^2+2ey+1}{\sqrt{2N}}\le\frac{11(1+y)^2}{\sqrt N}.
\end{equation}

\emph{Step 2.} Write $z=s_ke^t$, $w=2z$ and $R=R(w)$. By the formulas for
$I_0'$ and $I_1'$,
\[
 \beta_k'(t)=1+\frac{z\mathcal I_k'(z)}{\mathcal I_k(z)}
 =1+\frac{wR+c_kw-c_kR}{1+c_kR}=w+\frac{1+w(c_k-1)(1-R)}{1+c_kR}.
\]
At $t=\log u$ we have $w=w_k$, so $\beta_k'(\log u)=w_k+B_k$ with $B_k$ the
last fraction. By Lemma~\ref{lem:fisher-bessel}, $0<R<1$ and $w(1-R)\le1$.
So the numerator of $B_k$ lies in $[1,c_k]$ and is at most
$1+wc_k=1+y$. Hence $0<B_k\le1+y$. Subtracting $\frac12$ from
$1/(1+c_kR)\le B_k\le c_k/(1+c_kR)$ gives
\[
 -\frac{c_k-1}2\le B_k-\frac12\le\frac{(c_k-1)+c_k(1-R)}{2(1+c_kR)}
 \le\frac{c_k-1}2+\frac{c_k^2}{2y},
\]
by $c_k(1-R)\le c_k/w_k=c_k^2/y$. If $W^2\le\frac12$ on $K=k$, then
$c_k^2\le2$ and $c_k-1\le W^2$, since $(1-v)^{-1/2}-1$ is convex, vanishes
at $0$ and is below $\frac12$ at $v=\frac12$. So
$|B_k-\frac12|\le W^2/2+1/y$ there.

\emph{Step 3.} Let $\psi(a)=1-a^2/2-\sqrt{1-a^2}$ for $|a|\le1$.
The series of $\sqrt{1-v}$ is $1-v/2-\sum_{j\ge2}b_jv^j$ with $b_j\ge0$ and
$\sum_jb_j=\frac12$, so $0\le\psi(a)\le a^4/2$. On $K=k$,
$w_k=y(1-W^2/2)-y\psi(W)$. So on $0<K<N$ we have
$Z-G=(Z-\beta_K'(\log u))+(B_K-\frac12)-y\psi(W)$, and on the atoms
$K\in\{0,N\}$ we have $Z=0$, $W^2=1$ and $Z-G=-(y+1)/2$.

\emph{Step 4.} By Lemma~\ref{lem:fisher-W} and Markov's inequality,
$\E W^4\le8/T^2$, $\E W^8\le CT^{-4}$ and
$\Prob(W^2>\frac12)\le16\,\E W^8\le CT^{-4}$. Also
$\Prob(K\in\{0,N\})=p^{N-1}\le e^{-q(N-1)}\le e^{1/2-T}$.

\emph{Step 5.} By Step~3 and Minkowski's inequality, $\lVert Z-G\rVert_2$ is
at most the sum of the $L^2$ norms of the three terms of Step~3 on $0<K<N$ and of
$\frac{y+1}2$ on the atoms. By \eqref{eq:fisher-step1} the first is at most
$99T^2/\sqrt N$, as $(1+y)^2\le9T^2$. For the second we use Step~2 where
$W^2\le\frac12$ and $|B_K-\frac12|\le1+y$ elsewhere, so its square is at most
$\E W^4/2+2/y^2+C(1+y)^2T^{-4}\le C/T^2$. The third is at most
$\frac y2(\E W^8)^{1/2}\le C/T$, and the last is
$\frac{y+1}2\Prob(K\in\{0,N\})^{1/2}\le CTe^{-T/2}\le C/T$. So
$\E(Z-G)^2\le C(T^4/N+T^{-2})$, and dividing by $\Var J\ge T/4$ gives the
bound. Only Step~1 used $N\ge1700$, so
$N_0=1700$ works.
\end{proof}
